\documentclass[11pt]{article}
\usepackage[margin=1in]{geometry}
\title{The Orchard Below}
\author{Fictional Movie Catalog}
\date{2026}

\begin{document}
\maketitle

\section*{Catalog Metadata}
\textbf{Release date:} May 8, 2026\\
\textbf{Runtime:} 111 minutes\\
\textbf{Genres:} Mystery, Drama\\
\textbf{Director:} Ruben Hale\\
\textbf{Writer:} Sofia Bell\\
\textbf{Top cast:} Leona Price as Mara Bell; Milo Saye as Owen Bell;
Nora Wren as Elise Venn; Cal Danner as Thomas Bell.\\
\textbf{Mock IMDb rating:} 7.7 from 9,360 votes

\section*{Synopsis}
After their father's death, estranged siblings Mara and Owen Bell return to the
orchard where they grew up. A sinkhole opens beside the oldest pear tree and
reveals a sealed observatory beneath the roots. Inside are weather journals,
star charts, and recordings made by their mother, who vanished twenty years
earlier.

The recordings describe a rare atmospheric phenomenon that makes distant
sounds seem to come from beneath the ground. Mara believes the notebooks can
explain their mother's disappearance, while Owen suspects their father used the
story to hide a family crime. As an unseasonal frost threatens the harvest,
they reconstruct the observatory's last experiment and uncover a choice their
parents made to protect them.

\section*{Themes and Style}
The mystery is grounded in family memory rather than the supernatural. The
orchard changes with each discovery: bright spring rows gradually give way to
fog, bare branches, and the blue light of the underground dome. The film
focuses on grief, unreliable memory, forgiveness, and the difference between a
secret kept for control and one kept out of love.

\section*{Production Notes}
The orchard exteriors were filmed across four seasons and combined into a
single fictional week. The observatory set used a hand-operated brass
planetarium. The fictional production is rated PG for mature themes and mild
peril.

\end{document}
